% !TeX program = xelatex
% Компиляция: xelatex TiAl_600K_1273K.tex
% Исходник самодостаточен; данные таблиц включены в текст.
\documentclass[11pt,a4paper]{article}
\usepackage{fontspec}
\usepackage[russian]{babel}
\IfFontExistsTF{Cambria}{\setmainfont{Cambria}}{\setmainfont{DejaVu Serif}}
\IfFontExistsTF{Calibri}{\setsansfont{Calibri}}{\setsansfont{DejaVu Sans}}
\IfFontExistsTF{Consolas}{\setmonofont[Scale=MatchLowercase]{Consolas}}{\setmonofont[Scale=MatchLowercase]{DejaVu Sans Mono}}
\usepackage[a4paper,left=23mm,right=21mm,top=20mm,bottom=19mm,headheight=15pt,headsep=6mm,footskip=10mm]{geometry}
\usepackage{amsmath,amssymb,booktabs,array,xcolor,titlesec,fancyhdr,caption,hyperref}
\definecolor{ink}{HTML}{213542}
\definecolor{teal}{HTML}{006D70}
\definecolor{muted}{HTML}{526673}
\hypersetup{unicode=true,colorlinks=true,linkcolor=teal,urlcolor=teal,
 pdftitle={Равновесие фаз α и α₂ в Ti–Al при 600 и 1273,15 К},
 pdfsubject={Прямой MTP, эксперимент, кластерное разложение и CALPHAD}}
\urlstyle{same}
\captionsetup{font=small,labelfont={bf,color=teal},labelsep=period,singlelinecheck=false}
\titleformat{\section}{\sffamily\bfseries\Large\color{ink}}{\color{teal}\thesection}{.5em}{}
\titleformat{\subsection}{\sffamily\bfseries\large\color{ink}}{\thesubsection}{.5em}{}
\titlespacing*{\section}{0pt}{12pt plus 2pt}{5pt}
\titlespacing*{\subsection}{0pt}{9pt plus 2pt}{4pt}
\setlength{\parindent}{1.1em}
\setlength{\parskip}{2.4pt plus .6pt}
\setlength{\emergencystretch}{2em}
\brokenpenalty=10000 \widowpenalty=8000 \clubpenalty=8000
\renewcommand{\arraystretch}{1.16}
\raggedbottom
\pagestyle{fancy}\fancyhf{}
\fancyhead[L]{\small\sffamily\color{muted}Ti--Al\enspace/\enspace Равновесие ГПУ-фаз}
\fancyhead[R]{\small\sffamily\color{muted}600 и 1273,15 К}
\fancyfoot[R]{\small\sffamily\color{muted}\thepage}
\renewcommand{\headrulewidth}{.25pt}

\usepackage{pgfplots}
\pgfplotsset{compat=1.16}
\definecolor{linegray}{HTML}{C5D0D3}
\definecolor{rust}{HTML}{A8492B}
\begin{document}


\newcommand{\LowAlpha}{5{,}87}
\newcommand{\LowAlphaLo}{5{,}80}
\newcommand{\LowAlphaHi}{5{,}93}
\newcommand{\LowOrdered}{24{,}27}
\newcommand{\LowOrderedLo}{24{,}24}
\newcommand{\LowOrderedHi}{24{,}28}
\newcommand{\FixedAlpha}{7{,}57}
\newcommand{\FixedOrdered}{23{,}82}
\newcommand{\ChiMiddle}{0{,}529}
\newcommand{\ChiLarge}{0{,}542}

\thispagestyle{plain}
\noindent{\sffamily\small\bfseries\color{teal} НАУЧНЫЙ ОТЧЁТ
\hfill Ti--Al \(\cdot\) 25 ноября 2026 г.\par}
\vspace{3mm}
\noindent{\sffamily\bfseries\fontsize{23}{27}\selectfont\color{ink}
Границы \(\boldsymbol{\alpha/\alpha_2}\) в Ti--Al\\
при 600 и 1273,15 К\par}
\vspace{1.5mm}\noindent{\sffamily\large\color{muted}
Прямой расчёт на MTP при нулевом давлении, C46, CALPHAD и эксперимент\par}
\vspace{2mm}\noindent{\color{teal}\rule{\linewidth}{1pt}}\par

\noindent При \textbf{600 К} прямой классический расчёт на опубликованном MTP
даёт \(\boldsymbol{\LowAlpha}\) ат.\% Al в \(\alpha\) и
\(\boldsymbol{\LowOrdered}\) ат.\% Al в \(\alpha_2\). В основной
ячейке из 1400 атомов объём и отношение осей \(c/a\) независимы,
тепловые смещения ограничены условием \(W<0{,}40\). Условные 95\%-е
интервалы по временным блокам составляют
\([\LowAlphaLo;\,\LowAlphaHi]\) и \([\LowOrderedLo;\,\LowOrderedHi]\) ат.\% Al.

При \textbf{1273,15 К} измеренные границы составляют \(19{,}5\) и
\(23{,}0\) ат.\% Al. На этой изотерме ни прямой расчёт MTP, ни
Монте-Карло с C46 не разрешили два равновесных состава. Поэтому
численные границы MTP для этой температуры остаются неопределёнными.

\begin{center}\small
\begin{tabular}{@{}lrrr@{}}
\toprule
Модель или измерение & 600 К & 1000 К & 1273,15 К\\
\midrule
Прямой MTP, свободное \(c/a\) & \(\LowAlpha/\LowOrdered\) & 14,98/22,44 & две фазы не разрешены\\
C46, классический QHA--MC    & 4,55/24,10 & 14,90/21,92 & нет касательной\\
C46, квантовый QHA--MC       & 4,59/24,09 & 14,96/21,96 & \(\dagger\)\\
CALPHAD, Liang                & 10,93/20,04 & 14,75/21,00 & 19,35/23,19\\
Эксперимент, Kainuma и соавт. & --- & --- & 19,5/23,0\\
\bottomrule
\end{tabular}
\end{center}
\noindent{\small В ячейках приведены \(x_\alpha/x_{\alpha_2}\), ат.\% Al.
\(\dagger\) В исходном архиве квантового C46 обе границы разрешены при 1200~К,
при 1225, 1250 и 1300~К касательная с двумя различными контактами не найдена.
Прочерки в строке эксперимента означают отсутствие пары при точно этой
температуре в указанной таблице.}

\begin{center}\begin{tikzpicture}
\begin{axis}[width=.93\linewidth,height=72mm,xmin=565,xmax=1470,ymin=2,ymax=29,
  xlabel={\(T\), К},ylabel={Al, ат.\%},xtick={600,800,1000,1200,1400},
  xticklabels={600,800,1000,1200,1400},ytick={5,10,15,20,25},
  tick label style={font=\sffamily\footnotesize},label style={font=\sffamily\small},
  axis x line*=bottom,axis y line*=left,ymajorgrids,y grid style={linegray!65},
  legend style={font=\sffamily\footnotesize,draw=none,fill=white,at={(.5,-.27)},anchor=north,legend columns=2,column sep=5pt}]
\addplot[teal,line width=1.05pt,mark=*,mark size=1.4pt] coordinates {(600.000000,10.930179) (1000.000000,14.746957) (1273.150000,19.348045) (1323.150000,20.739924) (1373.150000,22.638277) (1423.150000,26.093040)};
\addlegendentry{CALPHAD, Liang}
\addplot[teal,line width=1.05pt,mark=*,mark size=1.4pt,forget plot] coordinates {(600.000000,20.037197) (1000.000000,21.004360) (1273.150000,23.194097) (1323.150000,23.957308) (1373.150000,25.055121) (1423.150000,27.219079)};
\addplot[muted,dashed,mark=triangle,mark size=1.9pt] coordinates {(600.000000,4.590143) (700.000000,7.210211) (800.000000,9.687890) (900.000000,12.154617) (1000.000000,14.958145) (1100.000000,18.425710) (1200.000000,23.464372)};
\addlegendentry{Квантовый C46, MC}
\addplot[muted,dashed,mark=triangle,mark size=1.9pt,forget plot] coordinates {(600.000000,24.089723) (700.000000,23.396689) (800.000000,22.661038) (900.000000,22.142876) (1000.000000,21.958271) (1100.000000,22.594336) (1200.000000,24.907841)};
\addplot[ink,only marks,mark=square,mark size=2pt,error bars/.cd,y dir=both,y explicit] coordinates {(1273.15,19.5) +- (0,0.5) (1323.15,21.3) +- (0,0.3) (1373.15,23.5) +- (0,0.2) (1423.15,24.7) +- (0,0.2)};
\addlegendentry{Эксперимент, Kainuma}
\addplot[ink,only marks,mark=square,mark size=2pt,forget plot,error bars/.cd,y dir=both,y explicit] coordinates {(1273.15,23.0) +- (0,0.3) (1323.15,23.5) +- (0,0.3) (1373.15,24.9) +- (0,0.2) (1423.15,25.7) +- (0,0.2)};
\addplot[rust,only marks,mark=diamond*,mark size=2.9pt] coordinates {(600.000000,5.865747) (600.000000,24.265025) (1000.000000,14.981744) (1000.000000,22.444830)};
\addlegendentry{Прямой MTP, свободное \(c/a\)}
\end{axis}\end{tikzpicture}
\captionof{figure}{Ti-богатые ветви сольвуса. Для CALPHAD и C46 линии соединяют
рассчитанные изотермы. Каждый ромб соответствует одной фазе при
температуре, где прямой расчёт разрешил две отдельные фазы. Планки у квадратов
передают точность из экспериментальной таблицы.}
\end{center}

\clearpage

\section*{При каких температурах измерены обе ветви}
В таблице~1 статьи Кайнумы, Пальма и Индена приведены составы на
Ti-богатой границе между разупорядоченным ГПУ-раствором \(\alpha\) и
упорядоченной фазой \(\alpha_2\) типа D0$_{19}$.
При 1000 и 1050~\(^\circ\)C авторы отжигали по 200~ч диффузионные пары
с исходными составами 15 и 25~ат.\% Al. Равновесные составы определяли
по концентрационным профилям у границы фаз электронно-зондовым
микроанализом. При 1100~\(^\circ\)C приведены точечные измерения в
двухфазном образце. Знак \(\pm\) в таблице передаёт указанную в
первоисточнике точность измерений.
\href{https://doi.org/10.1016/0966-9795(94)90018-3}{Kainuma, Palm, Inden (1994), табл.~1, с.~326.}

\begin{center}\small
\begin{tabular}{@{}rrrrrr@{}}
\toprule
&&\multicolumn{2}{c}{Эксперимент, ат.\% Al}&
\multicolumn{2}{c}{Liang, ат.\% Al}\\
\cmidrule(lr){3-4}\cmidrule(l){5-6}
\(T,\ ^\circ\)C & \(T,\) К & \(\alpha\) & \(\alpha_2\) & \(\alpha\) & \(\alpha_2\)\\
\midrule
1000 & 1273,15 & \(19{,}5\pm0{,}5\) & \(23{,}0\pm0{,}3\) & 19,35 & 23,19\\
1050 & 1323,15 & \(21{,}3\pm0{,}3\) & \(23{,}5\pm0{,}3\) & 20,74 & 23,96\\
1100 & 1373,15 & \(23{,}5\pm0{,}2\) & \(24{,}9\pm0{,}2\) & 22,64 & 25,06\\
1150 & 1423,15 & \(24{,}7\pm0{,}2\) & \(25{,}7\pm0{,}2\) & 26,09 & 27,22\\
\bottomrule
\end{tabular}
\end{center}
Две ближайшие к 1300~К строки соответствуют разным изотермам. При
1273,15~К эксперимент разрешает интервал шириной \(3{,}5\)~п.п., при
1323,15~К его ширина составляет \(2{,}2\)~п.п. Поэтому для прямого
атомистического теста выбрана первая температура. В таблице~1 независимой
работы Онумы и соавт. при 1373,15~К для \(\alpha/\alpha_2\) приведены
23,5 и 24,8~ат.\% Al. Их строка при 1273,15~К относится к
\(\alpha_2/\gamma\), то есть к границе с другой фазой.
\href{https://doi.org/10.1016/S1359-6454(00)00118-X}{Ohnuma et al. (2000), табл.~1.}

Для 600~К прямой пары составов в проверенных экспериментальных
таблицах и нанесённых в проекте точках Блэкберна не найдено.
Для этой изотермы количественный внешний ориентир даёт база Liang.
Во всех строках CALPHAD рассчитан заново при \(p=0\), с базой Liang
и десятью фазами бинарного сечения в \texttt{pycalphad}~0.11.2.
Отбиралась Ti-богатая ветвь с \(x_\alpha<x_{\alpha_2}\).
Для 600 и 1273,15~К расчёт только с \texttt{HCP\_A3} и
\texttt{ALTI3\_D019} дал те же контакты, что и расчёт с десятью фазами.
База и обработка фиксируют, с какой именно термодинамической оценкой
сопоставляется MTP. \href{https://doi.org/10.1007/s11669-024-01124-9}{Liang et al. (2024).}

\subsection*{Что даёт опыт при выборе эталона}
На изотерме 1273,15~К CALPHAD отличается от измерений Кайнумы на
\(-0{,}15\) п.п. для \(\alpha\) и \(+0{,}19\) п.п. для \(\alpha_2\).
Оба расхождения укладываются в приведённую авторами точность.
При 1323,15~К разности составляют \(-0{,}56\) и \(+0{,}46\) п.п.
Близость экспериментальных и расчётных границ зависит от температуры,
поэтому усреднение соседних точек вокруг 1300~К здесь скрыло бы часть
измеренного изменения.

На прежнем проектном рисунке точки Блэкберна оцифрованы с
рис.~13 Николса и соавт. Вблизи 1000~\(^\circ\)C они дают примерно
19 и 24~ат.\% Al. Для численного сравнения выше использована исходная
таблица Кайнумы. Николс и соавторы также исследовали переход на MTP
по скачку состава в двух сериях при заданных химических полях,
с начальными структурами \(\alpha\) и D0$_{19}$. Для MTP выше
1100~К петля гистерезиса стала слишком мала для определения границ
этим способом. Порог наблюдаемости петли зависит от метода и времени
наблюдения, поэтому при 1273,15~К мы отдельно проверили относительные
равновесные веса.
\href{https://doi.org/10.3390/met16020140}{Nichols et al. (2026), разд.~2.4, табл.~2, рис.~13.}

\clearpage

\section*{Что считается атомистическим равновесием}
В расчёте использован тот же опубликованный потенциал моментных тензоров
Qi и соавт., что и в серии при 1000~К. Состояние содержит типы всех
узлов, координаты атомов, объём и отношение осей гексагональной ячейки.
При заданном \(T,\ p=0\) в ходе Монте-Карло предлагаются замены
Ti\(\rightleftarrows\)Al, смещения атомов, изменения объёма и формы.
Химические степени свободы и ангармонические колебания входят в один
классический ансамбль. Для смены объёма учитывается множитель
\((V'/V)^N\). Один проход содержит \(N\) попыток химической замены и
\(N\) попыток смещения. Объём ограничен диапазоном 15--21~\AA$^3$/атом,
\(c/a\) в основном расчёте меняется в диапазоне 1,45--1,72.
В механическом контроле фиксируется авторское \(c/a=1{,}56374\);
объём продолжает флуктуировать. Каждой изотерме соответствуют ячейки
\(4\times4\times3\), \(6\times6\times4\), \(8\times8\times6\) и
\(10\times10\times7\) из двухатомных ячеек ГПУ.
\href{https://doi.org/10.1103/PhysRevMaterials.7.103602}{Qi et al. (2023).}

Исследуется равновесие кристаллических \(\alpha\) и \(\alpha_2\) на
ГПУ-решётке. Чтобы исключить ранее обнаруженное нефизическое сближение
Ti--Ti у исходного MTP, для каждого атома ограничена проекция смещения
от идеального узла вдоль каждой из 12 ближайших связей ГПУ. Величина
\(W\) --- максимум этих проекций, нормированных на половину длины соответствующей связи,
после удаления общего переноса. Основному расчёту соответствует
\(W<0{,}40\). Для объёмных и осевых проб дополнительно требуется
\(d_{\rm nn}^{\rm ideal}(V,c/a)(1-W_{\max})\ge1{,}68052\)~\AA,
где \(d_{\rm nn}^{\rm ideal}\) --- кратчайшее расстояние между идеальными узлами.
В рамках такого определения бассейна возможная конкуренция
со структурами вне ГПУ остаётся отдельной задачей.

Пусть \(n_j\) --- число атомов Al на каждой из четырёх подрешёток,
\(k=\sum_jn_j\). Параметр дальнего порядка
\[
 Q=\frac{4(4\max_j n_j-k)}{3N}
\]
равен единице для идеального Ti$_3$Al типа D0$_{19}$. В конечной
разупорядоченной ячейке \(Q\) положителен из-за операции максимума.
В окне \(l\) к энергии прибавлялся член
\(\frac12 K_l(Q-q_l)^2-\Delta\mu_l k\).
Химическое поле \(\Delta\mu_l\) относится к конфигурационной замене;
пересчёт к физическим массам даёт общий сдвиг поля при данной
температуре и сохраняет равновесные составы.

По измеренным \(k,Q\) метод MBAR снимает зонтики и восстанавливает
относительные веса состояний. При одинаковых \(N,T,W,p\), механическом
условии на \(c/a\) и допустимом пространстве конфигураций общая
атомистическая энергия сокращается в отношениях весов между окнами.
Это определение относительных свободных энергий позволяет обойтись
без абсолютной свободной энергии референсного кристалла. Для геометрически
допустимых проб изменение энергии вычисляется непосредственно по MTP. В каждой
изотерме используется собственное \(k_{\rm B}T\).
Когда две моды разделены минимумом \(Q^\ddagger\), порядок складывает
четыре эквивалентных варианта D0$_{19}$. В конечной ячейке оценка поля
сосуществования для одного варианта упорядочения определяется из уравнения
\[
 P_{\Delta\mu^*}(Q>Q^\ddagger)=4P_{\Delta\mu^*}(Q\le Q^\ddagger).
\]
По обе стороны разделителя вычисляют условные средние
\(\langle k/N\rangle\). Блочное ресэмплирование сохраняет временные
корреляции внутри траекторий. Получаемые интервалы условны по отношению
к размеру ячейки, определению бассейна и покрытию ансамбля окнами.
Если два отдельных бассейна не разрешены, два произвольных среза по \(Q\)
дали бы разные составы даже у одной фазы. В этой ситуации анализируют
изотерму состава, моды \(P(Q)\), восприимчивость и зависимость от размера.
\href{https://doi.org/10.1063/1.2978177}{Shirts, Chodera (2008).}

\subsection*{Граница применимости эталона}
На 288 узлах при 1273,15~К и поле 3,48~эВ уменьшение \(W_{\max}\)
с 0,40 до 0,35 повысило \(\langle Q\rangle\) с 0,29 до 0,40,
а среднюю концентрацию --- на 0,31 п.п. Таким образом, определение
кристаллического бассейна заметно для степени порядка. Соответствующие
профили и расчёт при фиксированном \(c/a\) проверены отдельно. Перенос
результата на неограниченный потенциал требует устранения нефизической
короткодействующей ветви и повторной проверки свободных энергий.

\clearpage

\section*{600 К. Две фазы и сходимость контактов}
Низкотемпературный переход даёт два хорошо разделённых максимума
распределения порядка. С увеличением ячейки барьер между ними растёт,
а условные средние составы в бассейнах меняются значительно слабее.
В квадратных скобках приведены 95\%-е блочные интервалы.
Столбец \(\beta\Delta F\) даёт высоту барьера в профиле
\(-\ln P(Q)\) при равных высотах пиков. Сам состав определён по
равенству статистических весов отдельных вариантов фаз.

\begin{center}\footnotesize\setlength{\tabcolsep}{4.6pt}
\begin{tabular}{@{}rrlccr@{}}\toprule
\(N\)& \(W_{\max}\) & \(c/a\) & \(x_\alpha\), ат.\% Al & \(x_{\alpha_2}\), ат.\% Al & \(\beta\Delta F\)\\\midrule
96 & \(0{,}40\) & своб. & \(5{,}93\ [5{,}72;\,6{,}18]\) & \(24{,}25\ [24{,}17;\,24{,}31]\) & \(4{,}9\)\\
288 & \(0{,}40\) & своб. & \(5{,}85\ [5{,}66;\,6{,}04]\) & \(24{,}35\ [24{,}31;\,24{,}37]\) & \(9{,}9\)\\
768 & \(0{,}40\) & своб. & \(5{,}87\ [5{,}81;\,5{,}96]\) & \(24{,}28\ [24{,}25;\,24{,}30]\) & \(21{,}2\)\\
1400 & \(0{,}40\) & своб. & \(5{,}87\ [5{,}80;\,5{,}93]\) & \(24{,}27\ [24{,}24;\,24{,}28]\) & \(35{,}7\)\\
288 & \(0{,}35\) & своб. & \(5{,}66\ [5{,}40;\,6{,}03]\) & \(24{,}24\ [24{,}19;\,24{,}28]\) & \(9{,}6\)\\
768 & \(0{,}35\) & своб. & \(5{,}87\ [5{,}77;\,5{,}98]\) & \(24{,}26\ [24{,}21;\,24{,}31]\) & \(21{,}7\)\\
1400 & \(0{,}35\) & своб. & \(5{,}87\ [5{,}79;\,5{,}98]\) & \(24{,}31\ [24{,}29;\,24{,}34]\) & \(36{,}9\)\\
288 & \(0{,}40\) & фикс. & \(7{,}52\ [7{,}43;\,7{,}62]\) & \(23{,}78\ [23{,}70;\,23{,}84]\) & \(6{,}4\)\\
768 & \(0{,}40\) & фикс. & \(7{,}57\ [7{,}49;\,7{,}68]\) & \(23{,}82\ [23{,}76;\,23{,}86]\) & \(12{,}2\)\\
\bottomrule\end{tabular}\end{center}

На 1400 узлах 32 зонтичных окна дали поле
\(\Delta\mu^*=3{,}156552\)~эВ, 95\%-й интервал [3,156251; 3,156886].
Средний порядок составляет 0,030 в \(\alpha\) и 0,970 в \(\alpha_2\),
барьер достиг \(35{,}7\,k_{\rm B}T\). Изменение контактов относительно
768 атомов составляет \(-0{,}005\) и \(-0{,}011\) п.п. Оценки числа
эффективных кадров с поправкой на корреляции --- 47 и 136 соответственно.
Восемь независимых траекторий без зонтика при полях 3,156 и 3,159~эВ
дали после приведения к \(\Delta\mu^*\) 5,81--5,83 и 24,30~ат.\% Al.
Разности с основным расчётом не превышают 0,06 п.п. Блочные интервалы
этих проверок и зонтичного расчёта перекрываются.

Ужесточение ограничения до \(W<0{,}35\) на 1400 атомах изменяет
\(x_\alpha\) на \(+0{,}006\)~п.п. (95\% от \(-0{,}104\) до \(+0{,}135\)),
а \(x_{\alpha_2}\) на \(+0{,}048\)~п.п. (от \(+0{,}017\) до \(+0{,}087\)).
Небольшая остаточная чувствительность правой границы различима
при достигнутой статистике. Эти две стенки характеризуют конечное
изменение ограничения в проверенном диапазоне.

\subsection*{Откуда берётся расхождение при 600 К}
В базе Liang ширина равновесного интервала составляет 9,11~п.п.,
в прямом MTP --- 18,40~п.п. Левая граница MTP ниже оценки Liang на 5,06~п.п.,
правая выше на 4,23~п.п. В упорядоченной фазе Al занимает 97,05\% мест
своей номинальной подрешётки. Для сравнения, при
1000~К прямой MTP и CALPHAD дают близкие левые границы
(14,98 и 14,75), тогда как правая граница прямого расчёта уже там
лежит на 1,44~п.п. выше значения Liang.

Когда форма ячейки свободна, \(\langle c/a\rangle\) при 600~К составляет
1,57764 в \(\alpha\) и 1,61012 в \(\alpha_2\). При фиксированном
авторском отношении осей на 768 атомах прямой MTP даёт
\(\FixedAlpha/\FixedOrdered\)~ат.\% Al. Для тех же механических условий
классический C46 даёт 4,55/24,10 на 5120 узлах. При согласованном
\(c/a\) прямой расчёт сдвигает левую границу относительно C46 на
\(+3{,}03\) п.п., правую --- на \(-0{,}29\) п.п.

При сравнении C46 с прямым MTP при фиксированном \(c/a\) обе модели
используют классическую статистику и один потенциал. Остаточная
разность включает аппроксимацию свободной энергии кластерным базисом,
квазигармоническое приближение и эффект ограничения атомных смещений
в прямом расчёте. Разделить эти вклады по двум равновесным составам
нельзя. Для оценки ошибки самого MTP при 600~К требуется независимый
ориентир, особенно для растворимости Al в \(\alpha\).

\clearpage

\section*{1273,15 К. Разрешаются ли две фазы?}
При \(N=96\) параметр порядка имеет два неглубоких максимума вблизи
\(\Delta\mu=3{,}46\)~эВ. Высота перемычки при равных высотах пиков
всего около \(0{,}64\,k_{\rm B}T\). На 288 узлах максимум и плечо в
этой области уже не дают двух равновесных бассейнов равной высоты.
В отличие от низкотемпературного результата, одно наличие плеча
в конечной ячейке не позволяет определить два предельных состава.

На 1400 узлах независимые траектории без зонтика, начатые из раствора
и упорядоченной структуры, при \(\Delta\mu=3{,}42\)~эВ дают
21,05 и 21,08~ат.\% Al, при 3,45~эВ --- 22,60 и 22,52~ат.\% Al.
В экспериментальном интервале 19,5--23,0~ат.\% Al максимум \(\chi_x\)
по перевзвешенной изотерме составляет \(\ChiMiddle\)~эВ\(^{-1}\) при
\(N=768\) и \(\ChiLarge\)~эВ\(^{-1}\) при \(N=1400\). В середине
интервала (21,25~ат.\% Al) средний глобальный порядок падает с 0,098
до 0,070 с увеличением ячейки от 768 до 1400 атомов.

Для проверки использована связь между откликом на химическое поле и
дисперсией состава, \(\chi_x=\partial\langle x\rangle/\partial\Delta\mu
=N\operatorname{var}(x)/(k_{\rm B}T)\). Если бы две фазы с разностью
составов 3,5~п.п. присутствовали с равными суммарными весами, один
межфазный вклад дал бы \(\chi_x=2{,}14\)~эВ\(^{-1}\) на 768 узлах
и 3,91 на 1400. Это оценка разрешимости при таком соотношении весов.
На 768 узлах флуктуации порядка максимальны около 27,62~ат.\% Al,
однако средние составы нижней и верхней четвертей распределения \(Q\)
там практически совпадают (27,59 и 27,63~ат.\%). На 1400 узлах при
3,56~эВ два старта дали 27,94 и 27,96~ат.\% Al, хотя средние \(Q\)
различались (0,090 и 0,263). Выше 30~ат.\% Al наибольшая ячейка
проверена прямыми точками около 30,4, 33,3 и 36,3~ат.\%. Перекрытие
между этими точками слабое. Узкий переход при
другом составе и характер потери дальнего порядка остаются открытыми.

В классическом C46 на 5120 узлах средние составы 51 окна покрывают
13,98--36,86~ат.\% Al. Поле возрастает между каждой парой соседних окон,
наименьший прирост составляет 6,53~мэВ. Во всех 384 блочных выборках
сохраняется монотонность. Внутренняя касательная с двумя контактами
не найдена ни в одной выборке. При перевзвешивании на квантовые
коэффициенты результат тот же. Эти данные дают независимую проверку
на уровне кластерной модели.

\begin{center}\begin{tikzpicture}
\begin{axis}[width=.92\linewidth,height=60mm,xmin=17.5,xmax=35,ymin=0,ymax=1.05,
  xlabel={Средний состав, ат.\% Al},ylabel={\(\langle Q\rangle\)},xtick={18,20,22,24,26,28,30,32,34},
  axis x line*=bottom,axis y line*=left,ymajorgrids,y grid style={linegray!65},
  tick label style={font=\sffamily\footnotesize},label style={font=\sffamily\small},
  legend style={font=\sffamily\footnotesize,draw=none,fill=white,at={(.98,.97)},anchor=north east}]
\path[fill=teal!8] (axis cs:19.5,0) rectangle (axis cs:23,1.03);
\addplot[muted,dotted,line width=1.1pt,no marks,unbounded coords=jump] coordinates {(17.018887,0.226560) (17.287697,0.233310) (17.558230,0.240481) (17.830540,0.248108) (18.104645,0.256225) (18.380513,0.264864) (18.658057,0.274053) (18.937124,0.283818) (19.217489,0.294176) (19.498852,0.305137) (19.780836,0.316702) (20.062987,0.328860) (20.344780,0.341587) (20.438532,0.345949) (20.625629,0.354844) (20.904893,0.368582) (21.181900,0.382734) (21.455957,0.397226) (21.726375,0.411969) (21.992491,0.426868) (22.253685,0.441822) (22.509402,0.456729) (22.759169,0.471487) (23.002609,0.485997) (23.239445,0.500167) (23.469511,0.513914) (23.692747,0.527163) (23.837795,0.535688) (23.909198,0.539852) (24.119006,0.551929) (24.322399,0.563354) (24.519681,0.574096) (24.711219,0.584138) (24.897433,0.593466) (25.078778,0.602078) (25.255742,0.609978) (25.428828,0.617172) (25.598549,0.623675) (25.765420,0.629500) (25.929955,0.634666) (26.092654,0.639190) (26.254011,0.643093) (26.414499,0.646394) (26.574577,0.649112) (26.734683,0.651267) (26.895233,0.652876) (27.056622,0.653957) (27.219222,0.654528) (27.383380,0.654604) (27.549414,0.654202) (27.717620,0.653338) (27.888259,0.652027) (28.061568,0.650286) (28.237747,0.648131) (28.416965,0.645578) (28.477404,0.644642) (28.599358,0.642646) (28.785025,0.639352) (28.974031,0.635717) (29.166408,0.631762) (29.362153,0.627508) (29.561234,0.622978) (29.763593,0.618194) (29.969147,0.613183) (30.177796,0.607967) (30.389430,0.602571) (30.603928,0.597020) (30.821173,0.591337) (31.041050,0.585544) (31.189044,0.581633) (31.263454,0.579665) (31.488293,0.573718) (31.715491,0.567724) (31.944987,0.561699) (32.176734,0.555660) (32.410700,0.549622) (32.646858,0.543596) (32.885183,0.537596) (33.125645,0.531632) (33.368202,0.525713) (33.612788,0.519850) (33.859308,0.514049) (34.107630,0.508320) (34.357574,0.502668)};
\addlegendentry{\(N=96\)}
\addplot[rust,dashed,line width=1.1pt,no marks,unbounded coords=jump] coordinates {(18.409079,0.137302) (18.612018,0.139377) (18.814030,0.141476) (19.015268,0.143614) (19.215856,0.145804) (19.415907,0.148065) (19.615527,0.150416) (19.814830,0.152879) (20.013930,0.155480) (20.212949,0.158246) (20.412005,0.161207) (20.611207,0.164396) (20.810649,0.167844) (21.010399,0.171584) (21.210498,0.175649) (21.410948,0.180068) (21.611714,0.184869) (21.812724,0.190075) (22.013863,0.195704) (22.214985,0.201769) (22.415910,0.208277) (22.616436,0.215227) (22.816341,0.222611) (23.015392,0.230414) (23.213358,0.238615) (23.410011,0.247184) (23.605142,0.256084) (23.798565,0.265273) (23.990121,0.274702) (24.179687,0.284316) (24.367177,0.294052) (24.552540,0.303846) (24.735764,0.313627) (24.916867,0.323322) (25.095900,0.332855) (25.272940,0.342151) (25.448084,0.351133) (25.621446,0.359730) (25.793155,0.367875) (25.963348,0.375504) (26.132172,0.382564) (26.299775,0.389009) (26.466316,0.394802) (26.631952,0.399918) (26.796852,0.404338) (26.961188,0.408057) (27.125143,0.411073) (27.288907,0.413397) (27.452680,0.415043) (27.616668,0.416031) (27.781081,0.416386) (27.946123,0.416133) (28.111985,0.415303) (28.278838,0.413927) (28.446818,0.412039) (28.616021,0.409674) (28.786495,0.406871) (28.958239,0.403670) (29.131203,0.400114) (29.305302,0.396249) (29.480419,0.392124) (29.656428,0.387787) (29.833205,0.383288) (30.010645,0.378676) (30.188669,0.373996) (30.367237,0.369291) (30.546338,0.364597) (30.725992,0.359945) (30.906239,0.355360) (31.087124,0.350862) (31.268688,0.346462) (31.450958,0.342169) (31.633938,0.337984) (31.817615,0.333906) (32.001955,0.329931) (32.186922,0.326052) (32.372485,0.322263) (32.558634,0.318553) (32.745397,0.314915) (32.932850,0.311339) (33.121124,0.307819)};
\addlegendentry{\(N=288\)}
\addplot[black!65,dashdotted,line width=1.1pt,no marks,unbounded coords=jump] coordinates {(17.132709,0.069841) (17.373861,0.070922) (17.612788,0.072040) (17.852948,0.073242) (18.097071,0.074564) (18.346570,0.076021) (18.385441,0.076257) (18.601271,0.077603) (18.859641,0.079276) (19.119497,0.081001) (19.378921,0.082749) (19.636906,0.084513) (19.893449,0.086310) (19.972184,0.086874) (20.149157,0.088174) (20.404686,0.090145) (20.653801,0.092200) (20.660358,0.092256) (20.916090,0.094527) (21.171521,0.096958) (21.426170,0.099534) (21.543292,0.100763) (21.679512,0.102225) (21.931005,0.105004) (22.180124,0.107856) (22.426470,0.110783) (22.613970,0.113098) (22.669885,0.113806) (22.910543,0.116955) (23.057485,0.118971) (23.148972,0.120265) (23.386007,0.123764) (23.622677,0.127480) (23.786860,0.130193) (23.860012,0.131438) (24.098822,0.135665) (24.339472,0.140187) (24.525736,0.143882) (24.581766,0.145028) (24.824995,0.150201) (25.068156,0.155703) (25.310285,0.161510) (25.550756,0.167572) (25.667244,0.170594) (25.789419,0.173805) (26.008343,0.179604) (26.026538,0.180085) (26.262567,0.186243) (26.497917,0.192070) (26.732870,0.197338) (26.967653,0.201815) (27.202595,0.205294) (27.438200,0.207605) (27.456368,0.207731) (27.553396,0.208268) (27.675030,0.208621) (27.913450,0.208267) (28.153412,0.206529) (28.394427,0.203475) (28.635737,0.199273) (28.876553,0.194185) (28.931980,0.192920) (29.116182,0.188558) (29.354035,0.182781) (29.589599,0.177237) (29.822508,0.172247) (30.052699,0.168035) (30.280542,0.164710) (30.367706,0.163670) (30.506830,0.162271) (30.732627,0.160632) (30.959086,0.159642) (31.187324,0.159117) (31.418343,0.158870) (31.652900,0.158744) (31.780794,0.158688) (31.891286,0.158640) (32.133074,0.158536) (32.377054,0.158473) (32.621543,0.158518) (32.865024,0.158717) (33.106803,0.159061) (33.273382,0.159348) (33.347324,0.159478) (33.587947,0.159851) (33.830300,0.160046) (34.075516,0.159941) (34.323732,0.159451) (34.574106,0.158543) (34.786698,0.157463) (34.825364,0.157238) (35.076527,0.155602) (35.327388,0.153732)};
\addlegendentry{\(N=768\)}
\addplot[teal,solid,line width=1.1pt,no marks,unbounded coords=jump] coordinates {(19.462204,0.062843) (19.684960,0.063529) (19.910222,0.064388) (20.138526,0.065360) (20.368628,0.066385) (20.597957,0.067415) (20.824175,0.068428) (21.046366,0.069446) (21.059316,0.069507) (21.265011,0.070532) (21.481133,0.071770) (21.695510,0.073238) (21.908413,0.074977) (22.119802,0.076981) (22.329637,0.079195) (22.538090,0.081532) (22.562541,0.081811) (22.745585,0.083891) (22.952696,0.086180) (23.160007,0.088338) (23.368023,0.090345) (23.577140,0.092212) (23.787571,0.093971) (23.999200,0.095654) (24.036636,0.095945) (24.211497,0.097285) (24.423649,0.098885) (24.535630,0.099729) (24.634878,0.100482) (24.844729,0.102108) (25.028692,0.103592) (25.053142,0.103795) (25.260315,0.105564) (25.466491,0.107419) (25.514879,0.107867) (25.671839,0.109349) (25.876460,0.111332) (26.080465,0.113348) (26.284021,0.115380) (26.475360,0.117290) (26.487314,0.117409) (26.690475,0.119410) (26.893593,0.121351) (27.096841,0.123190) (27.288582,0.124794) (27.300591,0.124889) (27.505387,0.126404) (27.711743,0.127681) (27.919876,0.128642) (28.129585,0.129187) (28.340316,0.129211) (28.551369,0.128648) (28.762088,0.127503) (28.910362,0.126389) (28.971999,0.125862) (29.180910,0.123859) (29.388946,0.121639) (29.596424,0.119321) (29.803551,0.117008) (30.010151,0.114799) (30.215672,0.112805) (30.371724,0.111494) (30.419514,0.111135) (30.621415,0.109866) (30.821632,0.109023) (31.020841,0.108580) (31.219947,0.108473) (31.419915,nan) (31.621626,nan) (31.825688,nan) (32.032212,nan) (32.240686,nan) (32.450039,nan) (32.658943,0.110621) (32.866300,0.110459) (33.071750,0.110107) (33.251897,0.109693) (33.275883,0.109634) (33.480006,0.109115) (33.685634,0.108600) (33.894124,0.108100) (34.106609,nan) (34.323918,nan) (34.546070,nan) (34.771578,nan) (34.997433,nan) (35.220264,nan) (35.437972,0.106785)};
\addlegendentry{\(N=1400\)}
\end{axis}\end{tikzpicture}
\captionof{figure}{Химический порядок при 1273,15 К на разных размерах,
\(W<0{,}40\), свободное \(c/a\). Тонированная полоса отмечает
экспериментальный интервал 19,5--23,0 ат.\% Al. Показаны точки с оценкой
эффективного числа кадров после поправки на временную корреляцию не менее 10.
Положительное среднее \(Q\) в конечной ячейке само по себе совместимо с
однофазным раствором.}
\end{center}

\clearpage

\section*{Что вносит кластерное разложение C46}
В C46 энергия статической решётки и колебательная свободная энергия
аппроксимированы по 46 исходным декорациям, 20 кластерным функциям и
девяти объёмам при фиксированном \(c/a=4{,}6046/2{,}9446\). Сохранены
авторские маски обучения и фононной устойчивости. Из 409 устойчивых пар
«структура--объём» 346 входят в подгонку, 63 служат для контроля.
Для каждого спектра классическая квазигармоническая метка вычисляется как
\[
 F_{\rm cl}(T)=\frac{T}{T_0}F_{\rm cl}(T_0)-gk_{\rm B}T\ln(T/T_0),
 \qquad T_0=1000~\text{К},\quad g=\sum_\nu w_\nu.
\]
Здесь \(w_\nu\) --- веса мод в расчёте на атом. Для каждого спектра \(g\)
и \(F_{\rm cl}(T_0)\) восстановлены в расчёте при 1000~К из авторских
таблиц квантовой колебательной свободной энергии. Формула следует из
суммы \(k_{\rm B}T\ln(\hbar\omega_\nu/k_{\rm B}T)\) при фиксированном объёме.
При каждом \(T\) коэффициенты разложения заново подгоняются к меткам.
В конфигурационном Монте-Карло коэффициенты интерполируются по объёму,
который релаксирует при \(p=0\). На контрольных декорациях среднеквадратичная
ошибка аппроксимации колебательной свободной энергии по отдельным объёмам
лежит в диапазоне 0,25--1,71~мэВ/атом при 600~К и
0,53--3,63~мэВ/атом при 1273,15~К.

Для интегрирования по составу VCSGC в каждом окне к энергии добавляется
\(N[k_{\rm B}T\kappa(x-c_0)^2-\mu_0 x]\), \(\kappa=100\). Среднее
значение \(x\) даёт оценку производной свободной энергии для конечной ячейки,
\(g'(\langle x\rangle)\simeq\mu_0-2k_{\rm B}T\kappa(\langle x\rangle-c_0)\).
Интегрирование кусочно-кубического
интерполянта PCHIP, сохраняющего локальную монотонность, даёт \(g(x)\)
с точностью до константы.
Для пары фаз ищутся два внутренних контакта общей опорной касательной
с разнесением не менее 0,5~п.п. На 5120 узлах использованы 52 окна при
600~К и 51 при 1273,15~К, по 3000 проходов уравновешивания и 9000 сбора.

При 600~К классическая модель на 768 и 5120 узлах даёт соответственно
4,53/24,15 и 4,55/24,10~ат.\% Al. На большем размере условные интервалы
составляют [4,49; 4,61] и [24,08; 24,13]. Перевзвешивание тех же
конфигураций на квантовые колебательные коэффициенты даёт 4,60/24,09,
независимый авторский квантовый расчёт --- 4,59/24,09. Для \(\alpha\)
изменение при замене классических колебаний квантовыми составляет
\(+0{,}050\) п.п. с 95\%-м интервалом [0,026; 0,078], для \(\alpha_2\) ---
\(-0{,}010\) п.п. с интервалом [\(-0{,}024\); 0,004].

\textbf{Химические корреляции.} В приближении Брегга--Вильямса
с теми же классическими коэффициентами общая касательная даёт
2,47/24,28 при 600~К и 15,44/22,05 при 1273,15~К. После перехода к
квантовым коэффициентам получаются соответственно 2,49/24,27 и
15,49/22,07. На верхней изотерме среднее поле предсказывает ширину
6,61~п.п. Сопоставление с VCSGC показывает, насколько вывод о двухфазном
равновесии чувствителен к учёту корреляций между заполнениями узлов.

\subsection*{Следующий расчёт}
Для определения температурного интервала замыкания сольвуса целесообразно
добавить прямые изотермы от 1100 до 1250~К. На каждой нужны профили
порядка и состава для нескольких размеров при одинаковой форме образца.
Рост межфазного барьера и два сохраняющихся состава позволят проследить
линию первого рода. Вблизи исчезновения барьера нужны более крупные
ячейки и независимые затравки. Для различения локального порядка и
структурного перехода вблизи 27~ат.\% Al следует сохранять временные ряды
структурного фактора концентрационных волн D0$_{19}$ и сравнить по размерам
кумулянт Биндера и отношение корреляционной длины к размеру ячейки. Независимая
проверка разбавленных растворов и антиструктурных дефектов Ti$_3$Al по DFT
позволит оценить вклад самого MTP в низкотемпературное расхождение.

\medskip\noindent{\small\textbf{Архив расчёта.} В 34 заданиях UMT на аккаунте
\texttt{u3129} завершены 563 траектории прямого MTP. После прогрева сохранено
189\,150 кадров при 600~К и 286\,560 при 1273,15~К. Все архивы сверены
по SHA-256. Исходник отчёта включает численные данные обоих рисунков.}

\end{document}

